A partir do desenvolvimento do sistema, foi possível aplicar os conhecimentos adquiridos ao longo do curso de Engenharia de Software, relacionando conceitos teóricos com a criação de uma solução voltada a uma necessidade real. O projeto permitiu compreender a importância da organização das informações e do acompanhamento das atividades para a rotina de um escritório de advocacia.

Os resultados obtidos mostram que a centralização dos dados em uma única plataforma contribui para facilitar o acesso às informações e melhorar o controle das demandas. Ao mesmo tempo, o trabalho evidenciou limitações que indicam caminhos para evolução: a ausência de integração com os sistemas dos tribunais, a inexistência de uma versão para dispositivos móveis, a falta de testes formais com usuários e a restrição dos relatórios gerenciais a indicadores básicos.

Como trabalhos futuros, sugere-se a integração com APIs de tribunais para automatizar a atualização de processos, o desenvolvimento de uma versão mobile, a realização de testes de usabilidade com usuários reais e a possibilidade de integrar recursos de inteligência artificial à plataforma, de modo a gerar relatórios gerenciais de forma automatizada e a apoiar a análise das informações cadastradas, oferecendo subsídios mais rápidos para a tomada de decisão no escritório.

Dessa forma, conclui-se que o JurisFlow representa uma proposta de solução tecnológica para apoiar a gestão de escritórios de advocacia, contribuindo para uma melhor organização das informações e das atividades do dia a dia.